%
% File: conclusion.tex
% Author: Diar Begisbayev
%
\chapter{Conclusion}
\label{chap:conclusion}
\setlength{\parskip}{1em}

We presented DTS-GSSF, a hybrid deep learning architecture for real-time passenger flow prediction on the Astana bus network. By integrating a Gated Recurrent Encoder for temporal feature extraction, adaptive graph propagation for spatial diffusion, and multi-head temporal attention for long-range dependency modelling, DTS-GSSF achieves an $R^2$ of 0.889 and a mean absolute error of 2.43 passengers per station-hour on a synthetic but realistic dataset of 1.3 million records across 374 stations. Ablation studies confirm that the TemporalAttention mechanism provides the largest single architectural improvement ($+1.3\%$ $R^2$), while the Negative Binomial likelihood appropriately captures the overdispersed nature of count-based ridership data.

Our key findings are threefold. First, gated recurrent encoders offer a compelling alternative to deeper recurrent and convolutional temporal backbones for spatio-temporal forecasting, providing stable gradient propagation and linear-time per-step processing. Second, combining a fixed physical adjacency matrix with a learnable adaptive matrix enables the model to respect known route topology while discovering latent passenger transfer patterns. Third, LoRA-based parameter-efficient adaptation makes it practical to specialize the model to individual routes without full retraining, reducing deployment cost in operational settings.

\section{Practical Implications}
\label{sec:practical_implications}

The results of this thesis carry direct implications for transit operators and dispatching services in Astana and comparable mid-sized cities. The model's inference latency of 4.2\,ms per sample means that demand forecasts for all 374 stations can be updated within approximately 1.6 seconds, satisfying the real-time requirements of dynamic vehicle dispatching systems. The Negative Binomial output provides not only point forecasts but also calibrated prediction intervals, enabling dispatchers to distinguish between routine demand fluctuations and unusually high boarding events that warrant proactive intervention such as adding extra vehicles or issuing passenger advisories.

The LoRA-based adaptation pathway is particularly relevant for operational deployment. Rather than maintaining separate models for each of the city's 40+ bus routes, a single base model can be specialised to individual routes by updating fewer than 6\% of total parameters, reducing storage requirements and enabling rapid redeloyment when route topologies change. This aligns with the incremental learning paradigm explored in prior work on adaptive transportation models \cite{mektepbayeva2025adaptive}, where online updates reduced retraining latency by 60\%.

From a data-engineering perspective, the feature engineering pipeline developed in this thesis---comprising cyclical temporal encodings, weather integration, and route-aware spatial features---can be directly integrated with the automated data-preparation framework described in \cite{yedilkhan2025intelligent}, which reduced manual preprocessing time by 75\%. Together, these components form a deployable end-to-end system from raw data ingestion to actionable demand forecasts.

\section{Answers to Research Questions}
\label{sec:rq_answers}

\textbf{RQ1 -- Temporal Modelling}: The gated recurrent encoder significantly outperforms deeper recurrent and convolutional alternatives. Compared to LSTM and GRU baselines ($R^2 \approx 0.81$), DTS-GSSF achieves $R^2 = 0.889$, a 7.2\% absolute improvement. The ablation study confirms that removing the TemporalAttention module (which complements the GRE) drops $R^2$ by 1.0\%. These results demonstrate that the gated recurrent encoder with linear-time per-step processing provides stable gradient propagation across the 72-hour context window. The practical significance of this finding is that transit operators can rely on a single unified temporal backbone rather than tuning the depth or hidden size of recurrent cells for each deployment scenario.

\textbf{RQ2 -- Spatial Modelling}: The dual-adjacency propagation mechanism improves prediction accuracy over fixed-graph approaches. Replacing the combined physical + adaptive matrix with the physical matrix alone reduces test $R^2$ from 0.889 to 0.886, while using only the adaptive matrix yields 0.884. The learnable mixing coefficient $\alpha$ converged to $\approx 0.58$ during training, confirming that both components contribute positively. This validates the hypothesis that latent transfer patterns---such as passengers switching between routes at interchange stops---carry predictive information beyond what route topology alone provides.

\textbf{RQ3 -- Uncertainty Quantification}: The Negative Binomial likelihood provides calibrated uncertainty estimates for overdispersed count data. The learned dispersion parameter $\kappa = 7.2$ indicates moderate overdispersion (variance = $\mu + \mu^2/7.2$), confirming that the data violate the Poisson assumption. For operational dispatching, this means the model can produce prediction intervals that appropriately widen during peak hours when boarding variability is highest, rather than producing overconfident narrow intervals that a Gaussian likelihood would generate.

\textbf{RQ4 -- Computational Efficiency}: DTS-GSSF achieves competitive accuracy while maintaining linear-time per-step processing. Inference latency is 4.2\,ms per sample (238 predictions/second), well within the operational requirement for real-time bus dispatching. The 470K-parameter model is compact enough for edge deployment on embedded hardware typically available on transit vehicles, and the LoRA adaptation reduces trainable parameters by 94\%, making route-level specialisation feasible even on resource-constrained devices.

\textbf{RQ5 -- Transferability}: LoRA-based parameter-efficient adaptation enables rapid route-level specialization. By updating only the low-rank matrices, the number of trainable parameters is reduced by 94\% compared to full fine-tuning. In pilot experiments, route-specific adaptation converges in 12--18 epochs while retaining 98\% of the full-model test $R^2$. Across prediction horizons, the model maintains consistent accuracy: $R^2$ ranges from 0.874 at 15 minutes to 0.901 at 120 minutes, with central business district stations showing the highest predictability and peripheral residential stations exhibiting the largest residual error, likely due to sporadic boarding patterns at low-traffic stops.

\section{Limitations and Future Work}
\label{sec:limitations}

Several limitations remain. The model was evaluated on synthetic data generated from OpenStreetMap topology and heuristic demand simulation; performance on real-world ridership logs from the Astana transit operator may differ, particularly during anomalous events such as severe weather, public holidays, or service disruptions that the synthetic generator does not fully capture. The $R^2$ ceiling near $89\%$ suggests that a portion of variance is irreducible given the current feature set, likely reflecting unobserved special events, service disruptions, and stochastic individual behaviour. Finally, the temporal attention complexity of $O(N T^2 d)$ may become a bottleneck for networks requiring very long contexts, though linear-attention approximations can address this.

Promising directions for future work include: (i) extending the model to multi-modal transit networks by incorporating metro and ride-hailing data, following the cross-city federated approach explored in \cite{sakhipov2026deep}; (ii) replacing the handcrafted adaptive adjacency with a fully learned dynamic graph that evolves over time, building on the drift-aware mechanisms proposed in \cite{sakhipov2025federated}; (iii) exploring causal intervention modelling to predict the effect of service changes such as route cancellations or frequency adjustments; and (iv) deploying the model on edge devices for real-time inference at bus stops, leveraging the compact 470K-parameter architecture and LoRA adaptation pathway. We will release the training code, model weights, and synthetic dataset upon acceptance to facilitate reproducibility and community benchmarking.